\documentclass[11pt]{article}
\newcommand{\LabNumber}{7}
\newcommand{\LabTitle}{Recursion}
\newcommand{\LabTopic}{Base cases, recursive traces, binary search, Fibonacci}
\input{common.tex}

\begin{document}
\labheader

\section*{Objectives}
\begin{itemize}[leftmargin=*,nosep]
  \item Identify the base case(s) and recursive case(s) of a recursive method.
  \item Draw a recursion trace to understand how the call stack unwinds.
  \item Implement and analyse: factorial, binary search, power, and Fibonacci.
  \item Recognise the exponential cost of naive Fibonacci and apply memoisation.
\end{itemize}

\begin{summary}[Quick Reference]
\textbf{Anatomy of a recursive method.}\\
1.~\textbf{Base case} --- returns without a recursive call.\\
2.~\textbf{Recursive case} --- calls itself with a \emph{strictly smaller} sub-problem.

\textbf{Factorial.}
\begin{lstlisting}
static long factorial(int n) {
    if (n == 0) return 1;
    return n * factorial(n - 1);
}
\end{lstlisting}

\textbf{Binary search (recursive).}
\begin{lstlisting}
static int binSearch(int[] a, int lo, int hi, int target) {
    if (lo > hi) return -1;
    int mid = (lo + hi) / 2;
    if (a[mid] == target) return mid;
    else if (target < a[mid]) return binSearch(a, lo, mid-1, target);
    else                      return binSearch(a, mid+1, hi, target);
}
\end{lstlisting}

\textbf{Complexity.}\quad Factorial: $O(n)$. Binary search: $O(\log n)$. Naive Fibonacci: $O(2^n)$.
\end{summary}

\subsection*{Quick Experiments (Try It)}

\begin{tryit}{Stack overflow --- how deep can you go?}
Call \texttt{factorial(n)} for increasing \texttt{n} (try 10\,000, 100\,000).
At what value does a \texttt{StackOverflowError} occur?
\end{tryit}

\begin{tryit}{Counting recursive calls}
Add a static counter \texttt{int calls = 0} to your Fibonacci and increment at each call.
Compare for \texttt{fib(10)}, \texttt{fib(20)}, \texttt{fib(30)}. How does the count grow?
\end{tryit}

\section*{In-Lab Exercises (target: 2 hours)}

\begin{labexercise}{Factorial and Power \hfill {\small \itshape (20 min)}}
Implement:
\begin{itemize}[leftmargin=*,nosep]
  \item \texttt{long factorial(int n)} --- recursive.
  \item \texttt{long power(long base, int exp)} --- recursive; \texttt{power(b,0) = 1}.
\end{itemize}
Test: \texttt{factorial(10) = 3628800}, \texttt{power(2,10) = 1024}.
Draw the recursion trace for \texttt{power(3,4)}.
\end{labexercise}

\begin{labexercise}{Recursive Binary Search \hfill {\small \itshape (20 min)}}
Implement \texttt{int binSearch(int[] a, int lo, int hi, int target)}.
Test on a sorted array of 15 elements: search for the middle, first, last, and a missing element.
Count recursive calls for each search.
\end{labexercise}

\begin{labexercise}{Fibonacci: Naive vs.\ Memoised \hfill {\small \itshape (40 min)}}
\begin{enumerate}[leftmargin=*,nosep]
  \item Implement naive \texttt{long fibNaive(int n)}.
  \item Implement memoised \texttt{long fibMemo(int n, long[] memo)}: initialise \texttt{memo} to $-1$;
        store each result before returning.
  \item Time both for $n = 10, 20, 30, 40$ using \texttt{System.nanoTime()}.
\end{enumerate}
\begin{center}
\begin{tabular}{@{}rrrr@{}}
\toprule
$n$ & \texttt{fibNaive (ns)} & \texttt{fibMemo (ns)} & Speedup \\
\midrule
10 & & & \\ 20 & & & \\ 30 & & & \\ 40 & & & \\
\bottomrule
\end{tabular}
\end{center}
\end{labexercise}

\begin{labexercise}{Sum of Digits and Digit Count \hfill {\small \itshape (25 min)}}
Write two recursive methods (base case: $n < 10$; recursive case uses \texttt{n / 10}):
\begin{itemize}[leftmargin=*,nosep]
  \item \texttt{int digitSum(int n)} --- sum of all digits.
  \item \texttt{int countDigits(int n)} --- number of digits.
\end{itemize}
Test: \texttt{digitSum(1234) = 10}, \texttt{countDigits(100) = 3}.
\end{labexercise}

\begin{aicollab}
\textbf{Suggested prompts:}
\begin{itemize}[leftmargin=*,nosep]
  \item ``Draw the recursion tree for \texttt{fib(5)}, counting duplicate sub-problems.''
  \item ``Explain memoisation: how does it change Fibonacci from $O(2^n)$ to $O(n)$?''
\end{itemize}
\medskip\textbf{Verify:} Have the AI generate memoised Fibonacci. Run it alongside yours for $n=40$.
Do they agree? Is the AI's version iterative or recursive?
\end{aicollab}

\takehomebanner

\begin{trace}[Recursion trace for \texttt{power}]
Draw the full recursion trace (call stack) for \texttt{power(2, 5)}.
Show each call with its arguments and the return value as the stack unwinds.
\end{trace}

\begin{trace}[Binary search trace]
Trace \texttt{binSearch(a, 0, 7, 19)} on $a = [3, 7, 11, 15, 19, 23, 27, 31]$.
For each call show: \texttt{lo}, \texttt{hi}, \texttt{mid}, and the comparison result.
How many calls are made?
\end{trace}

\begin{writecode}[Recursive string reverse]
Write \texttt{String reverse(String s)} recursively.
Base case: \texttt{s} is empty or one character.
Recursive case: last character + \texttt{reverse(s without last character)}.
\end{writecode}

\vspace{4pt}
\noindent\textbf{Submission.} Save files as \texttt{lab7\_ex*.java} and submit on the course site.

\end{document}
